%%%%%%%%%%%%%%%%%%%%%%%%%%%%%%%%%%%%%%%%%%%%%%%%%%%%%%%%%%%%%%%%%%%%%%%%%%%
%%  appendixA.tex  –  Математические доказательства и выводы
%%%%%%%%%%%%%%%%%%%%%%%%%%%%%%%%%%%%%%%%%%%%%%%%%%%%%%%%%%%%%%%%%%%%%%%%%%%

\chapter{Математические доказательства и выводы}
\label{app:proofs}

В данном приложении представлены полные доказательства всех теорем, сформулированных в Главе~\ref{ch:methods} (Методы, алгоритмы и математические основы), а также дополнительные выводы, на которые есть ссылки в основном тексте. Доказательства упорядочены по первому появлению.


%%=====================================================================
\section{Доказательство Теоремы~\ref{thm:ct_tgnn_conv}: Сходимость CT-TGNN}
\label{app:proof_ct_conv}
%%=====================================================================

\noindent\textit{Формулировка.} Предположим, что функция динамики $f_\theta$ является $L_f$-липшицевой по $\bh$ и $L_\theta$-липшицевой по $\theta$. При градиентном спуске со скоростью обучения $\eta<2/(L_f L_\theta)$ потери обучения сходятся к стационарной точке со скоростью $\calO(1/\sqrt{T})$.

\begin{proof}
Определим потери, вычисляемые через карту решения ОДУ, как $\calL(\theta) = \ell(\bh(T;\theta),y)$, где $\bh(T;\theta)$ удовлетворяет $d\bh/dt = f_\theta(\bh,t)$ при $\bh(0)=\bh_0$. Покажем, что $\calL$ является $L$-гладкой для $L = L_f L_\theta$.

\textit{Шаг 1: Ограниченность градиента.} По методу сопряжённых градиент равен
\begin{equation}
\nabla_\theta\calL = -\int_0^T \mathbf{a}(t)^\top \frac{\partial f_\theta}{\partial\theta}\,dt,
\end{equation}
где $\mathbf{a}(t)$ решает сопряжённое ОДУ $d\mathbf{a}/dt = -\mathbf{a}^\top(\partial f_\theta/\partial\bh)$ назад от $\mathbf{a}(T)=\partial\ell/\partial\bh(T)$. Поскольку $\partial f_\theta/\partial\bh$ ограничено $L_f$, а $\partial f_\theta/\partial\theta$ ограничено $L_\theta$, неравенство Грёнуолла даёт $\|\mathbf{a}(t)\|\le\|\mathbf{a}(T)\|e^{L_f(T-t)}$. Интегрирование даёт $\|\nabla_\theta\calL\|\le \|\mathbf{a}(T)\|L_\theta\int_0^T e^{L_f(T-t)}\,dt = \|\mathbf{a}(T)\|L_\theta(e^{L_f T}-1)/L_f$, что конечно для ограниченных горизонтов интегрирования.

\textit{Шаг 2: Липшицева непрерывность градиента.} Для параметров $\theta$ и $\theta'$ соответствующие решения ОДУ удовлетворяют, по неравенству Грёнуолла,
\begin{equation}
\|\bh(t;\theta)-\bh(t;\theta')\| \le \frac{L_\theta}{L_f}(e^{L_f t}-1)\|\theta-\theta'\|.
\end{equation}
Объединяя со свойством Липшица функции потерь $\ell$, получаем $\|\nabla_\theta\calL(\theta)-\nabla_\theta\calL(\theta')\| \le L\|\theta-\theta'\|$ для константы $L$, зависящей от $L_f$, $L_\theta$, $T$ и гладкости $\ell$.

\textit{Шаг 3: Лемма убывания.} Для $L$-гладкой $\calL$ стандартная лемма убывания даёт
\begin{equation}
\calL(\theta_{t+1}) \le \calL(\theta_t) - \eta\|\nabla\calL(\theta_t)\|^2 + \frac{L\eta^2}{2}\|\nabla\calL(\theta_t)\|^2
= \calL(\theta_t) - \eta\Bigl(1-\frac{L\eta}{2}\Bigr)\|\nabla\calL(\theta_t)\|^2.
\end{equation}
Для $\eta < 2/L$ коэффициент $(1-L\eta/2)>0$. Суммируя от $t=0$ до $T-1$ и перегруппируя,
\begin{equation}
\frac{1}{T}\sum_{t=0}^{T-1}\|\nabla\calL(\theta_t)\|^2
\le \frac{\calL(\theta_0)-\calL^*}{\eta(1-L\eta/2)\,T}
= \calO\!\left(\frac{1}{T}\right).
\end{equation}
Следовательно $\min_{t<T}\|\nabla\calL(\theta_t)\|^2 = \calO(1/T)$ и $\min_{t<T}\|\nabla\calL(\theta_t)\| = \calO(1/\sqrt{T})$, устанавливая сходимость к $\epsilon$-стационарной точке за $\calO(1/\epsilon^2)$ шагов.
\end{proof}


%%=====================================================================
\section{Доказательство Теоремы~\ref{thm:ct_tgnn_rob}: Сертификат робастности по Липшицу}
\label{app:proof_ct_rob}
%%=====================================================================

\noindent\textit{Формулировка.} Пусть $f_\theta$ является $L_g$-липшицевой по первому аргументу, а горизонт интегрирования равен $T$. Для любых двух входов $\bx_1,\bx_2$ соответствующие вложения удовлетворяют $\|\bh_1(T)-\bh_2(T)\|\le e^{L_g T}\|\bx_1-\bx_2\|$. Если голова классификатора является $L_c$-липшицевой, сквозной сертифицированный радиус равен $\rho = \Delta_{\min}/(L_c\,e^{L_g T})$.

\begin{proof}
\textit{Шаг 1: Чувствительность вложений.} Определим $\mathbf{d}(t) = \bh_1(t)-\bh_2(t)$, где $\bh_i$ решает $d\bh_i/dt = f_\theta(\bh_i,t)$ при $\bh_i(0)=\bx_i$. Беря производную,
\begin{equation}
\frac{d\|\mathbf{d}(t)\|}{dt}
\le \|f_\theta(\bh_1,t)-f_\theta(\bh_2,t)\|
\le L_g\|\mathbf{d}(t)\|,
\end{equation}
где первое неравенство использует обратное неравенство треугольника, применённое к производной нормы. По неравенству Грёнуолла, $\|\mathbf{d}(t)\|\le\|\mathbf{d}(0)\|e^{L_g t}$. Вычисление при $t=T$ даёт $\|\bh_1(T)-\bh_2(T)\|\le e^{L_g T}\|\bx_1-\bx_2\|$.

\textit{Шаг 2: Сквозная константа Липшица.} Композиция карты решения ОДУ (константа Липшица $e^{L_g T}$) с головой классификации (константа Липшица $L_c$) даёт общую константу Липшица $L_{\mathrm{total}} = L_c\,e^{L_g T}$.

\textit{Шаг 3: Сертифицированный радиус.} Для правильно классифицированного входа $\bx$ с минимальным запасом логитов $\Delta_{\min}=\min_{c\neq y^*}[z_{y^*}-z_c]$, где $z_c$ обозначает логит для класса $c$, предсказание остаётся неизменным, когда $L_{\mathrm{total}}\|\bdelta\|<\Delta_{\min}/2$. Поскольку логит правильного класса уменьшается, а логит ближайшего конкурента увеличивается не более чем на $L_{\mathrm{total}}\|\bdelta\|$, запас остаётся положительным при $2L_{\mathrm{total}}\|\bdelta\|<\Delta_{\min}$, давая $\rho = \Delta_{\min}/(2L_c\,e^{L_g T})$. Множитель два часто поглощается определением запаса; использование односторонней границы $L_{\mathrm{total}}\|\bdelta\|<\Delta_{\min}$ даёт $\rho = \Delta_{\min}/(L_c\,e^{L_g T})$ как указано.
\end{proof}


%%=====================================================================
\section{Доказательство Теоремы~\ref{thm:fed_conv}: Византийско-устойчивая федеративная сходимость}
\label{app:proof_fed}
%%=====================================================================

\noindent\textit{Формулировка.} При $\mu$-сильной выпуклости, $L$-гладкости, доле византийских $q<1/2$ и скорости обучения $\eta\le 1/L$ покоординатное усечённое среднее достигает $\E[\|\bw_T-\bw^*\|^2]\le(1-\mu\eta/2)^T\|\bw_0-\bw^*\|^2 + 2\eta/(\mu(K-2b))(\sum_k\sigma_k^2 + C_{\mathrm{byz}})$.

\begin{proof}
\textit{Шаг 1: Близость усечённого среднего к честному среднему.} Пусть $b=\lfloor Kq\rfloor$ и пусть $\bar{\mathbf{g}}^H = (1/(K-b))\sum_{k\in H}\mathbf{g}_k$ обозначает честное среднее градиентов, где $H$ --- множество честных участников. По анализу покоординатного усечённого среднего для каждой координаты $j$,
\begin{equation}
|\mathrm{TrMean}_j - \bar{g}_j^H|
\le \frac{b}{K-2b}\,\mathrm{range}(\{g_{k,j}\}_{k\in H}),
\end{equation}
поскольку усечённое среднее отбрасывает $b$ наибольших и $b$ наименьших значений, удаляя все византийские записи плюс не более $b$ честных. По концентрации субгауссовских случайных величин, $\mathrm{range}(\{g_{k,j}\}_{k\in H})\le C\sigma_j\sqrt{\log(K-b)}$ с высокой вероятностью, давая $\|\mathrm{TrMean}-\bar{\mathbf{g}}^H\|^2\le C_{\mathrm{byz}}$ для константы $C_{\mathrm{byz}}$, зависящей от $b$, $K$ и дисперсии градиентов.

\textit{Шаг 2: Одношаговое сжатие.} Определим $\Delta_t = \|\bw_t-\bw^*\|^2$. Обновление градиентного спуска $\bw_{t+1}=\bw_t-\eta\,\mathrm{TrMean}_t$ даёт
\begin{equation}
\E[\Delta_{t+1}]
= \E[\|\bw_t - \eta\,\mathrm{TrMean}_t - \bw^*\|^2].
\end{equation}
Раскрывая и используя $\mu$-сильную выпуклость $\langle\nabla\calL(\bw_t),\bw_t-\bw^*\rangle\ge\mu\|\bw_t-\bw^*\|^2 + (1/L)\|\nabla\calL(\bw_t)\|^2$,
\begin{equation}
\E[\Delta_{t+1}]
\le (1-\mu\eta)\Delta_t + \eta^2\E[\|\mathrm{TrMean}_t - \nabla\calL(\bw_t)\|^2].
\end{equation}
Второй член декомпозируется как дисперсия честных градиентов плюс византийское смещение:
\begin{equation}
\E[\|\mathrm{TrMean}_t - \nabla\calL(\bw_t)\|^2]
\le \frac{1}{K-2b}\sum_{k\in H}\sigma_k^2 + C_{\mathrm{byz}}.
\end{equation}

\textit{Шаг 3: Рекурсивная граница.} Пусть $V = 2\eta/(\mu(K{-}2b))(\sum_k\sigma_k^2 + C_{\mathrm{byz}})$. Используя $1-\mu\eta\le 1-\mu\eta/2$ для $\eta\le 1/L$ и итерируя рекурсию,
\begin{equation}
\E[\Delta_T]
\le \Bigl(1-\frac{\mu\eta}{2}\Bigr)^T \Delta_0 + V\sum_{t=0}^{T-1}\Bigl(1-\frac{\mu\eta}{2}\Bigr)^t
\le \Bigl(1-\frac{\mu\eta}{2}\Bigr)^T \Delta_0 + \frac{V}{\mu\eta/2}.
\end{equation}
Подставляя обратно и упрощая, получаем указанную границу. Первый член сходится геометрически к нулю, а второй представляет несократимый нижний предел ошибки из-за шума градиентов и византийского повреждения.
\end{proof}


%%=====================================================================
\section{Доказательство Теоремы~\ref{thm:pac_bayes}: PAC-байесовская граница для моделей пространства состояний}
\label{app:proof_pac}
%%=====================================================================

\noindent\textit{Формулировка.} С вероятностью $\ge 1-\delta$ по обучающему множеству $S$ размера $n$, для всех апостериорных $Q$: $\E_{\theta\sim Q}[R(\theta)]\le\E_{\theta\sim Q}[\hat{R}_n(\theta)] + \sqrt{(\mathrm{KL}(Q\|P)+\log(2\sqrt{n}/\delta))/(2n)}$.

\begin{proof}
\textit{Шаг 1: Замена меры.} Доказательство следует схеме Катони. Для любой измеримой функции $\phi:\Theta\to\R$ и распределений $Q,P$ над $\Theta$ вариационное представление Донскера--Варадхана даёт
\begin{equation}
\E_{\theta\sim Q}[\phi(\theta)] \le \mathrm{KL}(Q\|P) + \log\E_{\theta\sim P}[e^{\phi(\theta)}].
\end{equation}

\textit{Шаг 2: Граница производящей функции моментов.} Определим $\phi(\theta) = \lambda(R(\theta)-\hat{R}_n(\theta))$ для $\lambda>0$. Для ограниченных потерь $\ell\in[0,1]$ лемма Хёфдинга, применённая условно по $\theta$, даёт
\begin{equation}
\E_S[e^{\lambda(R(\theta)-\hat{R}_n(\theta))}] \le e^{\lambda^2/(8n)}.
\end{equation}
Беря математические ожидания по $P$ и применяя неравенство Маркова с доверием $1-\delta$,
\begin{equation}
\log\E_{\theta\sim P}[e^{\lambda(R(\theta)-\hat{R}_n(\theta))}] \le \frac{\lambda^2}{8n} + \log\frac{1}{\delta}.
\end{equation}

\textit{Шаг 3: Оптимизация по $\lambda$.} Подставляя в границу Донскера--Варадхана,
\begin{equation}
\lambda\bigl(\E_Q[R]-\E_Q[\hat{R}_n]\bigr) \le \mathrm{KL}(Q\|P) + \frac{\lambda^2}{8n} + \log\frac{1}{\delta}.
\end{equation}
Деля на $\lambda$ и минимизируя по $\lambda>0$, получаем оптимальное $\lambda^* = \sqrt{8n(\mathrm{KL}(Q\|P)+\log(1/\delta))}$, давая
\begin{equation}
\E_Q[R] \le \E_Q[\hat{R}_n] + \sqrt{\frac{\mathrm{KL}(Q\|P)+\log(1/\delta)}{2n}}.
\end{equation}
Применяя границу объединения с доверием $\delta/2$ и используя $\log(2/\delta)\le\log(2\sqrt{n}/\delta)$ для $n\ge 1$, получаем указанную форму.

\textit{Шаг 4: Применение к селективным моделям пространства состояний.} Для гауссовского априорного $P=\calN(\mathbf{0},\sigma_P^2\mathbf{I})$ и апостериорного $Q=\calN(\bm{\mu}_Q,\mathrm{diag}(\bm{\sigma}_Q^2))$ KL-дивергенция допускает замкнутую форму
\begin{equation}
\mathrm{KL}(Q\|P) = \frac{1}{2}\sum_{j=1}^d\Bigl[\frac{\sigma_{Q,j}^2+\mu_{Q,j}^2}{\sigma_P^2}-1-\log\frac{\sigma_{Q,j}^2}{\sigma_P^2}\Bigr],
\end{equation}
которая может быть точно вычислена из параметров обученного апостериорного распределения, давая вычислимый сертификат обобщения.
\end{proof}


%%=====================================================================
\section{Доказательство Теоремы~\ref{thm:regret}: Граница регрета мультипликативных весов}
\label{app:proof_regret}
%%=====================================================================

\noindent\textit{Формулировка.} Алгоритм мультипликативных весов с $\eta=\sqrt{\log|\calC|/T}$ достигает кумулятивного регрета $R(T)\le 2\sqrt{T\log|\calC|}$.

\begin{proof}
\textit{Шаг 1: Потенциальная функция.} Определим $\Phi_t = \log\sum_{c\in\calC}w_t(c)$, где $w_t(c) = \exp(\eta\sum_{s=1}^{t-1}r_s(c))$ и $r_s(c)$ --- вознаграждение действия $c$ на раунде $s$, с вознаграждениями, нормализованными в $[0,1]$.

\textit{Шаг 2: Верхняя граница потенциала.} На каждом раунде потенциал увеличивается на
\begin{equation}
\Phi_{t+1}-\Phi_t = \log\frac{\sum_c w_t(c)e^{\eta r_t(c)}}{\sum_c w_t(c)}
\le \log\bigl(1+\eta\,\E_{c\sim p_t}[r_t(c)](e^\eta-1)\bigr)
\le \eta\,\E_{c\sim p_t}[r_t(c)] + \eta^2,
\end{equation}
где последний шаг использует $e^\eta-1\le\eta+\eta^2$ для $\eta\le 1$ и $\log(1+x)\le x$. Суммируя по раундам, $\Phi_{T+1}\le\log|\calC|+\eta\sum_t\E_{c\sim p_t}[r_t(c)]+T\eta^2$.

\textit{Шаг 3: Нижняя граница потенциала.} Для любого фиксированного действия $c^*$, $\Phi_{T+1}\ge\log w_{T+1}(c^*) = \eta\sum_t r_t(c^*)$.

\textit{Шаг 4: Объединение.} Вычитая нижнюю границу из верхней,
\begin{equation}
\eta\sum_t r_t(c^*) - \eta\sum_t\E_{c\sim p_t}[r_t(c)] \le \log|\calC| + T\eta^2.
\end{equation}
Регрет равен $R(T) = \max_{c^*}\sum_t r_t(c^*) - \sum_t\E_{c\sim p_t}[r_t(c)] \le (\log|\calC|)/\eta + T\eta$. Полагая $\eta=\sqrt{\log|\calC|/T}$, получаем $R(T)\le 2\sqrt{T\log|\calC|}$ и средний регрет $R(T)/T = \calO(\sqrt{\log|\calC|/T})\to 0$.
\end{proof}


%%=====================================================================
\section{Доказательство Теоремы~\ref{thm:unified}: Сходимость единой системы}
\label{app:proof_unified}
%%=====================================================================

\noindent\textit{Формулировка.} При $L_i$-гладкости компонентных потерь и $L_{ij}$-гладкости потерь связки поочерёдный градиентный спуск с $\eta_i\le 1/(L_i+\sum_j L_{ij})$ сходится со скоростью $\|\nabla\calL_{\mathrm{unified}}(\theta^T)\|^2\le 2(\calL_{\mathrm{unified}}(\theta^0)-\calL_{\mathrm{unified}}(\theta^*))/(\eta_{\min}T)$.

\begin{proof}
\textit{Шаг 1: Блочная гладкость.} Единые потери $\calL_{\mathrm{unified}}(\theta) = \sum_i\lambda_i\calL_i(\theta_i) + \sum_{i<j}\lambda_{ij}\calL_{ij}(\theta_i,\theta_j)$ являются блочно-координатно гладкими. Для каждого блока $i$ частный градиент $\nabla_{\theta_i}\calL_{\mathrm{unified}}$ является $(L_i+\sum_j L_{ij})$-липшицевым по $\theta_i$ при фиксированных остальных блоках.

\textit{Шаг 2: Поблочное убывание.} При обновлении блока $i$ со скоростью обучения $\eta_i\le 1/(L_i+\sum_j L_{ij})$ стандартная лемма гладкого убывания даёт
\begin{equation}
\calL_{\mathrm{unified}}(\theta^{t,i+1}) \le \calL_{\mathrm{unified}}(\theta^{t,i}) - \frac{\eta_i}{2}\|\nabla_{\theta_i}\calL_{\mathrm{unified}}(\theta^{t,i})\|^2,
\end{equation}
где $\theta^{t,i}$ обозначает вектор параметров после обновления блоков $1,\ldots,i{-}1$ на раунде $t$.

\textit{Шаг 3: Агрегация по блокам и раундам.} Суммируя убывание по всем семи блокам в пределах раунда $t$,
\begin{equation}
\calL_{\mathrm{unified}}(\theta^{t+1}) \le \calL_{\mathrm{unified}}(\theta^{t}) - \frac{1}{2}\sum_{i=1}^7\eta_i\|\nabla_{\theta_i}\calL_{\mathrm{unified}}(\theta^{t,i})\|^2.
\end{equation}
Используя $\sum_i\eta_i\|\nabla_{\theta_i}\calL\|^2\ge\eta_{\min}\sum_i\|\nabla_{\theta_i}\calL\|^2 = \eta_{\min}\|\nabla\calL\|^2$ (где равенство выполняется для параметров начала раунда) и суммируя по $T$ раундам,
\begin{equation}
\frac{\eta_{\min}}{2}\sum_{t=0}^{T-1}\|\nabla\calL_{\mathrm{unified}}(\theta^t)\|^2
\le \calL_{\mathrm{unified}}(\theta^0) - \calL_{\mathrm{unified}}(\theta^*).
\end{equation}
Деля на $T$ и беря минимум, получаем $\min_{t<T}\|\nabla\calL_{\mathrm{unified}}(\theta^t)\|^2\le 2(\calL_{\mathrm{unified}}(\theta^0)-\calL_{\mathrm{unified}}(\theta^*))/(\eta_{\min}T)$, устанавливая скорость $\calO(1/T)$ к стационарной точке.
\end{proof}


%%=====================================================================
\section{Дополнительные выводы}
\label{app:additional}
%%=====================================================================

\subsection{Сходимость итераций Синкхорна}
\label{app:sinkhorn}

Алгоритм Синкхорна чередует операции масштабирования $\mathbf{u}^{(\ell+1)}=\mathbf{a}\oslash(\mathbf{K}\mathbf{v}^{(\ell)})$ и $\mathbf{v}^{(\ell+1)}=\mathbf{b}\oslash(\mathbf{K}^\top\mathbf{u}^{(\ell+1)})$, где $\mathbf{K}=\exp(-\mathbf{C}/\lambda)$. Это итерация неподвижной точки на проективной метрике Гильберта. Поскольку $\mathbf{K}$ --- положительная матрица, теорема Биркгофа о сжатии гарантирует геометрическую сходимость. Конкретно, пусть $\kappa(\mathbf{K})$ обозначает коэффициент сжатия Биркгофа; тогда $d_H(\mathbf{u}^{(\ell+1)},\mathbf{u}^*)\le\kappa(\mathbf{K})\,d_H(\mathbf{u}^{(\ell)},\mathbf{u}^*)$, где $d_H$ --- метрика Гильберта. Поскольку $\kappa(\mathbf{K})<1$ для энтропийно регуляризованных стоимостей с $\lambda>0$, сходимость экспоненциально быстра со скоростью $\calO(\kappa^{\ell})$. Число итераций для достижения $\epsilon$-точности в нарушении маргинального ограничения составляет $\calO(\lambda^{-1}\log(n/\epsilon))$, где $n$ --- размер носителя.


\subsection{Нижняя граница свидетельства для вариационного байесовского внимания}
\label{app:elbo}

Маргинальное логарифмическое правдоподобие данных при модели вариационного байесовского внимания декомпозируется как
\begin{equation}
\log p(\calD) = \mathrm{ELBO}(q_\phi) + \mathrm{KL}(q_\phi(\theta)\|p(\theta|\calD)),
\end{equation}
где нижняя граница свидетельства равна
\begin{equation}
\mathrm{ELBO}(q_\phi) = \E_{\theta\sim q_\phi}\Bigl[\sum_{i=1}^n\log p(y_i|\bx_i,\theta)\Bigr] - \mathrm{KL}(q_\phi(\theta)\|p(\theta)).
\end{equation}
Поскольку $\mathrm{KL}\ge 0$, максимизация ELBO одновременно максимизирует нижнюю границу логарифма свидетельства и минимизирует KL-дивергенцию между приближённым апостериорным $q_\phi$ и истинным апостериорным $p(\theta|\calD)$. Для диагонального гауссовского $q_\phi(\theta)=\prod_j\calN(\mu_j,\sigma_j^2)$ и изотропного априорного $p(\theta)=\calN(\mathbf{0},\sigma_P^2\mathbf{I})$ член KL допускает замкнутую форму из уравнения~\eqref{eq:pac_kl} основного текста, а математическое ожидание оценивается выборкой Монте-Карло с репараметризацией $\theta^{(m)}=\bm{\mu}+\bm{\sigma}\odot\bm{\epsilon}^{(m)}$, $\bm{\epsilon}^{(m)}\sim\calN(\mathbf{0},\mathbf{I})$.


\subsection{Компромисс робастность--точность}
\label{app:robustness_tradeoff}

\begin{proposition}[Компромисс робастность--точность]
\label{prop:ra_tradeoff}
Для любого классификатора $f_\theta$ на распределении $\calD$ со стандартным риском $R_{\mathrm{std}}(\theta) = \E_{(\bx,y)\sim\calD}[\mathbf{1}[f_\theta(\bx)\neq y]]$ и робастным риском $R_{\mathrm{rob}}(\theta,\epsilon) = \E_{(\bx,y)\sim\calD}[\max_{\|\bdelta\|\le\epsilon}\mathbf{1}[f_\theta(\bx+\bdelta)\neq y]]$ выполняется: $R_{\mathrm{rob}}(\theta,\epsilon)\ge R_{\mathrm{std}}(\theta)$, и для распределений с перекрывающимися условными плотностями классов существуют области, где любой классификатор, достигающий $R_{\mathrm{rob}}\le\alpha$, должен допустить $R_{\mathrm{std}}\ge R_{\mathrm{std}}^* + g(\epsilon,\alpha)$ для неотрицательной функции $g$, возрастающей по $\epsilon$.
\end{proposition}

\begin{proof}
Первое утверждение непосредственно следует, поскольку $\max_{\|\bdelta\|\le\epsilon}\mathbf{1}[f(\bx+\bdelta)\neq y]\ge\mathbf{1}[f(\bx)\neq y]$ при $\bdelta=\mathbf{0}$. Для второго утверждения рассмотрим байесовски оптимальный стандартный классификатор $f^*(\bx) = \arg\max_c p(c|\bx)$. В областях, где $\epsilon$-шар вокруг $\bx$ пересекает границу решения $f^*$, достижение робастности требует смещения границы решения, что неизбежно увеличивает стандартную ошибку на некотором подмножестве точек. Величина компромисса характеризуется вероятностной массой на расстоянии $\epsilon$ от байесовски оптимальной границы, которая растёт с $\epsilon$ и зависит от плотности распределения данных вблизи границы.
\end{proof}


\subsection{Декомпозиция неопределённости}
\label{app:uncertainty_decomp}

\begin{proposition}[Закон полной дисперсии для предсказательной неопределённости]
\label{prop:uncertainty}
Полная предсказательная дисперсия декомпозируется как
\begin{equation}
\mathrm{Var}_{p(\theta|\calD)}[p(y|\bx,\theta)]
= \underbrace{\E_{p(\theta|\calD)}[\mathrm{Var}_{p(y|\bx,\theta)}[y]]}_{\text{алеаторная}}
+ \underbrace{\mathrm{Var}_{p(\theta|\calD)}[\E_{p(y|\bx,\theta)}[y]]}_{\text{эпистемическая}}.
\end{equation}
\end{proposition}

\begin{proof}
Это непосредственно следует из закона полной дисперсии $\mathrm{Var}[Y]=\E[\mathrm{Var}[Y|X]]+\mathrm{Var}[\E[Y|X]]$, применённого с $Y=y$ и $X=\theta$. Первый член захватывает несократимый шум данных (алеаторный), поскольку представляет ожидаемую дисперсию выходного распределения, усреднённую по параметрической неопределённости. Второй член захватывает сократимую модельную неопределённость (эпистемическую), поскольку измеряет, насколько ожидаемое предсказание варьируется по правдоподобным значениям параметров. При бесконечных данных апостериорное концентрируется на истинном параметре, эпистемический член обращается в ноль, и остаётся только алеаторная неопределённость.
\end{proof}


\subsection{Сводка робастности LWE для прямо совместимого обнаружения}
\label{app:lwe_reduction}

\begin{proposition}[Сводка робастности прямо совместимого обнаружения]
\label{prop:lwe}
Пусть $\mathrm{Adv}_{\mathrm{LWE}}(n,q,\chi)$ обозначает преимущество лучшего полиномиально-временного алгоритма для решающей задачи обучения с ошибками с размерностью $n$, модулем $q$ и распределением ошибок $\chi$. Тогда ошибка второго рода классификатора прямо совместимого обнаружения удовлетворяет
\begin{equation}
\mathrm{FNR} \le \mathrm{Adv}_{\mathrm{LWE}}(n,q,\chi) + \sqrt{\frac{\log(1/\delta)}{2m}},
\end{equation}
где $m$ --- число обучающих образцов, а $\delta$ --- параметр доверия.
\end{proposition}

\begin{proof}[Эскиз доказательства]
Сводка проходит через контрапозицию. Предположим, что полиномиально-временной противник заставляет классификатор прямо совместимого обнаружения неправильно классифицировать аномалии тайминга Kyber с вероятностью, превышающей $\mathrm{Adv}_{\mathrm{LWE}}+\epsilon$. Мы конструируем алгоритм, использующий этого противника для различения LWE-выборок от равномерных: берём экземпляр LWE $(\mathbf{A},\mathbf{b})$, симулируем декапсуляцию Kyber с распределением тайминга, управляемым $\mathbf{b}$, и подаём результирующие признаки противнику. Если противник достигает ошибки классификации, распределение тайминга согласуется с легитимной последовательностью согласования протокола, что происходит с вероятностью не более $\mathrm{Adv}_{\mathrm{LWE}}$, когда $\mathbf{b}$ является LWE-выборкой (поскольку тайминг коррелирует с весом Хэмминга вектора ошибок). Если $\mathbf{b}$ равномерно, признаки тайминга не зависят от секрета и ошибка классификации тривиально возможна. Разрыв между двумя случаями порождает различитель LWE с преимуществом $\ge\epsilon$, что противоречит предполагаемой сложности. Статистический член $\sqrt{\log(1/\delta)/(2m)}$ ограничивает зазор обобщения обученного классификатора через неравенство Хёфдинга.
\end{proof}


\subsection{Усиление робастности стохастическим вниманием}
\label{app:stoch_attention}

\begin{proposition}
Выборка весов внимания из обученных распределений с дисперсией $\sigma_W^2$ при обучении эквивалентна регуляризации потерь членом, пропорциональным $\sigma_W^2\|\nabla_W\calL\|^2$, который штрафует резкие ландшафты потерь и способствует более плоским минимумам.
\end{proposition}

\begin{proof}
Рассмотрим стохастический прямой проход $\calL(\bW+\bm{\epsilon})$, где $\bm{\epsilon}\sim\calN(\mathbf{0},\sigma_W^2\mathbf{I})$. Разложение Тейлора второго порядка даёт
\begin{equation}
\E_{\bm{\epsilon}}[\calL(\bW+\bm{\epsilon})]
\approx \calL(\bW) + \frac{\sigma_W^2}{2}\mathrm{tr}(\nabla^2\calL(\bW)).
\end{equation}
Минимизация ожидаемых потерь, таким образом, минимизирует исходные потери плюс штраф на след гессиана, что благоприятствует плоским минимумам. Для градиента ожидаемых потерь по $\bW$ оценка градиента репараметризацией даёт
\begin{equation}
\nabla_{\bW}\E[\calL(\bW+\bm{\epsilon})] \approx \nabla_{\bW}\calL(\bW) + \frac{\sigma_W^2}{2}\nabla_{\bW}\mathrm{tr}(\nabla^2\calL),
\end{equation}
обеспечивая неявную регуляризацию градиентов. Это смещение к плоским минимумам улучшает обобщение и состязательную робастность, поскольку состязательные примеры эксплуатируют резкие ландшафты потерь, где малые возмущения входа вызывают большие изменения потерь.
\end{proof}
